% !PRINTED 412
\weberpage{412}
Wenn nun $\mu < n$ ist, so kann es nicht zugleich durch $a$ und durch $b$ theilbar sein,
und also können auch $a$ und $b$ nicht beide die kleinsten positiven Exponenten der
Potenzen von $\alpha$ und $\beta$ sein, die gleich 1 werden, d. h. also, wenn $r$ nicht
primitive $n^{\mathrm{te}}$ Einheitswurzel ist, so sind auch $\alpha$ und $\beta$ nicht
zugleich primitive $a^{\mathrm{te}}$ und $b^{\mathrm{te}}$ Einheitswurzeln, oder wenn
$\alpha$ und $\beta$ primitive $a^{\mathrm{te}}$ und $b^{\mathrm{te}}$ Einheitswurzeln
sind, so ist ihr Product $r$ primitive $n^{\mathrm{te}}$ Einheitswurzel. Auf der anderen
Seite ist klar, dass, wenn $\alpha$ oder $\beta$ zugleich Einheitswurzel von niedrigerem
Grade als $a$ oder $b$ ist, auch $r$ Einheitswurzel von niedrigerem als dem
$n^{\mathrm{ten}}$ Grade sein wird.

Zerf\"allt $n$ in mehrere Factoren $a, b, c \ldots$, von denen je zwei zu einander
relativ prim sind, und sind $\alpha, \beta, \gamma \ldots$ Einheitswurzeln der Grade
$a, b, c \ldots$, und setzt man

\weq{3}{r = \alpha\beta\gamma \ldots,}
so schliesst man durch mehrmalige Anwendung der vorigen S\"atze, dass in (3) alle
$n^{\mathrm{ten}}$ Einheitswurzeln, und jede nur einmal, enthalten sind, und ferner, dass
$r$ dann und nur dann \gesperrt{primitive} $n^{\mathrm{te}}$ Einheitswurzel ist, wenn
$\alpha, \beta, \gamma \ldots$ primitive Einheitswurzeln der Grade $a, b, c \ldots$ sind.

Bezeichnen wir jetzt die Anzahl der primitiven $n^{\mathrm{ten}}$ Einheitswurzeln durch
$\varphi(n)$, so folgt aus dem hier Bewiesenen

\weq{4}{\varphi(n) = \varphi(a)\varphi(b)\varphi(c) \ldots,}
wenn

\[
n = abc \ldots
\]
und $a, b, c \ldots$ Zahlen sind, die, je zwei und zwei, zu einander relativ prim sind.

Nun kann man jede Zahl $n$ auf eine und nur auf eine Weise in ein Product von
Primzahlpotenzen zerlegen

\[
n = p^\pi p_1^{\pi_1} p_2^{\pi_2} \ldots,
\]
worin $p, p_1, p_2 \ldots$ verschiedene Primzahlen und $\pi, \pi_1, \pi_2 \ldots$
positive Exponenten sind, so dass aus der Formel (4) folgt:

\weq{5}{\varphi(n) = \varphi(p^\pi)\, \varphi(p_1^{\pi_1})\, \varphi(p_2^{\pi_2}) \ldots,}
und dass es also nur noch darauf ankommt, zu entscheiden, ob und wie viele primitive
Einheitswurzeln des Grades $p^\pi$ existiren.

Diese Frage ist aber sehr einfach zu entscheiden. Wenn nämlich $\varrho$ eine
Einheitswurzel vom Grade $p^\pi$ ist, so ist der niedrigste Grad, zu dem $\varrho$ als
Einheitswurzel geh\"ort, ein Theiler von $p^\pi$, also eine Potenz von $p$, und wenn er
also nicht gleich
%CONT>>
